% ----------------------------------------------
% Capítol 1
% ----------------------------------------------
\chapter{David i la Resistència del LaTeX}

\epigraf{És més fàcil desintegrar un àtom que un prejudici.}{Albert Einstein}

\lettrine[loversize=0.2, lines=2]{M}{entre} el regnat de l'Eric, L'Homogeneïtzador, s'estenia amb fermesa i la pau governava l'univers dels arxius, una nova força va emergir des de les trinxeres del primer curs. Un jove estudiant anomenat David va començar de manera humil, resolent els críptics problemes de cinemàtica i penjant les seves solucions. Però la veritable revolució va arribar quan va tenir una revelació divina: el descobriment del LaTeX.

Allò que abans eren fotografies desenfocades de guixots fets amb pressa es va transformar sota les mans d'en David en documents d'una bellesa incalculable. Les equacions fluïen amb una precisió tipogràfica celestial i els teoremes es presentaven amb una elegància absoluta. Fascinat, en David va intentar predicar la paraula d'aquesta nova fe, difonent la religió del LaTeX entre els seus companys per purificar el OneDrive.

Tanmateix, els inicis van ser durs i la resposta, hostil. La comunitat va rebutjar la seva crida. Per a la immensa majoria d'estudiants, ofegats per les pràctiques i l'estrès, escriure codi resultava un martiri innecessari. Ningú volia aprendre comandaments estranys ni barallar-se amb errors de compilació quan fer una simple foto a una llibreta era instantani. El LaTeX era repudiat per ser considerat massa lent i molest, una eina reservada per a puristes incompresos.

Però en David no es va rendir. Ignorant el menyspreu i la mandra generalitzada, va continuar la seva croada solitària. Amb un fervor incansable, va transcriure, ordenar i omplir cada racó de la carpeta de primer, demostrant de manera pràctica que la llegibilitat i perfecció del format compensaven qualsevol esforç. El primer curs es va anar convertint en un santuari intocable de netedat acadèmica.

Des de les altes esferes de l'administració, l'Eric ho observava tot. L'Homogeneïtzador sentia una barreja d'admiració i d'una profunda por. En la seva ment de guardià, la irrupció d'aquest jove obstinat el feia dubtar. I si aquesta tirania del format dividia els estudiants i marginava aquells que només podien aportar escrits a mà? Durant mesos, l'Eric es va mantenir expectant, vigilant des de les ombres cada fitxer compilat que s'afegia al servidor.

Finalment, incapaç de negar que aquell estàndard era l'evolució definitiva, l'Eric va apropar-se. En un acte silenciós de reconeixement, en David va ser acollit oficialment a la família OneDrive i designat com a l'Hereu legítim per la seva empenta. No obstant això, el líder va ser dictaminant: el jove era prometedor, però massa inexperimentat. Abans de cedir-li el comandament i deixar-lo imposar la seva religió a la resta, l'hereu hauria de madurar i sobreviure als implacables camps de batalla dels cursos superiors.



% ----------------------------------------------
% Capítol 6
% ----------------------------------------------
\chapter{L'Alçament del Consell de Savis}

\epigraf{El govern és massa important per deixar-lo en mans d’un sol home.}{Winston Churchill}

\lettrine[loversize=0.2, lines=2]{E}{n} arribar la fi de l'aventura de l'Eric per la Universitat, va aparèixer un nou obstacle per a la dinastia OneDrive. L'innocent repositori, el qual havia estat una dòcil criatura en els seus inicis, s'havia tornat una fera, aprenent en el procés que el sagrat ordre establert que la mantenia lligada podia ser una arma de doble tall. La dedicació i poder necessaris per a dominar tal nombre d'arxius, PDF's i resums s'havien tornat excessius per a un sol mortal. Ni tan sols la promesa del jove David i la seva creixent revolució del LaTeX podien sostenir el pes de tot l'imperi.

Amb l'esperança de no perdre la raó en el camí, l'Eric OneDrive pregava a Newton, Dirac, Maxwell, Bagan, entre altres déus del panteó dels físics; per una solució. En sentir les seves súpliques i gran devoció, ells el van recompensar amb la gran revelació. Calia realitzar un canvi de règim, trencar amb l'absolutisme de la famosa dinastia i donar pas a una assemblea de ferms guerrers capaços de donar-ho tot per defensar el seu líder i guiar el jove hereu: el Consell de Savis; una unió entre els acadèmics més preparats de cada promoció per a recopilar els coneixements any rere any.

% ----------------------------------------------
% Capítol 7
% ----------------------------------------------
\chapter{OneDrive Thiago, L'Historiador}

\epigraf{“Els pobles que obliden la seva història estan condemnats a repetir-la.”}{George Santayana}

\lettrine[loversize=0.2, lines=2]{Q}{uan} va arribar l'hora, l'Eric va marxar triomfant, deixant enrere uns estudiants beneïts pel material col·lectiu i un Consell de Savis ja establert. Però al capdavant, mentre en David continuava la seva croada tipogràfica des de les bases, s'erigia un digne veterà com a paladí de les equacions per liderar la transició: en Thiago OneDrive. Ell va refinar l'estructura per adequar-la a totes les reestructuracions que el nou pla formatiu designava, un programa atrevit però no per això menys capriciós o críptic.

En Thiago, més enllà de desenvolupar una inèdita tasca de reorganització que esgarrifaria a la majoria, va fixar la mirada en expandir nous horitzons, cap a les cròniques perdudes dels ancestres. Soterrades entre les dunes del temps, en Thiago va descobrir dues ruïnes oblidades, intents dels nostres predecessors de centralitzar la saviesa del moment en reservoris de manuscrits ancestrals, desgraciadament sense èxit: el Google Drive Física i el Mega Física; tanmateix, l'últim d'aquests estava massa malmès per a poder rescatar cap document. Aquesta troballa converteix el OneDrive Física en la màxima culminació de dècades de llàgrimes, sacrifici i fisicaires perduts en la ignorància.

Després d'inestimables hores de treball, en Thiago va poder estendre els confins del OneDrive tant en el passat, present i futur; per prolongar el suport indispensable pels que tenen i tindran la voluntat de desxifrar els grans enigmes del cosmos, des de la quàntica fins a les galàxies més llunyanes.





















